\chapter{Radices and decimal arithmetic}
\label{ch:conv}

Numbers are stored in base $2^{64}$ but read and written in base 10. Because 10 is not a power
of two, converting is real arithmetic: every output digit depends on every input limb. Done
the obvious way it costs $n^2$; BigMath does it with a divide-and-conquer scheme that costs
a logarithmic factor more than one multiplication. The code is in \file{src/common/Parser.cpp}.

\section{Why the obvious method is slow}

The schoolbook way to parse is ``for each digit: $r \leftarrow 10r + \text{digit}$''. With $L$
digits that is $L$ big-number operations, each on a number of up to $L/19$ limbs: about
$L^2/19$ limb operations. Printing the schoolbook way (``divide by 10, emit the remainder'') is
equally quadratic. At a million digits that is around $10^{10}$ operations.

The library removes this cost in two layers: a constant-factor speedup (handle 18--19 digits
at a time) and an asymptotic one (divide and conquer).

\section{Small numbers: 18 or 19 digits at a time}
\label{sec:linparse}

\textbf{Parsing.} Instead of one digit, take 18 at a time: $10^{18} < 2^{63}$, so 18 digits fit in
one limb with headroom. Each step is $r \leftarrow r \cdot 10^{18} + \text{chunk}$: one
multiply-by-one-limb and one add over the current value. That is 18 times fewer big-number
operations than digit by digit.

\begin{example}[: parsing a 25-digit number]
\code{"1234567890123456789012345"} has 25 digits. $25 = 7 + 18$, so the first chunk is the
leading 7 digits and then one full 18-digit chunk:
\begin{align*}
r &= 1\,234\,567\\
r &= 1\,234\,567 \cdot 10^{18} + 890\,123\,456\,789\,012\,345 .
\end{align*}
Two steps instead of 25.
\end{example}

\textbf{Printing.} Divide by $10^{19}$ (the largest power of ten that fits in a limb) repeatedly;
each remainder is a block of 19 digits. Each division is the fast single-limb division of
\cref{sec:shortdiv}.

\begin{example}[: printing $2^{64}$]
$2^{64}$ is the limbs $[0, 1]$. Dividing by $10^{19}$ gives quotient $1$ and remainder
$8\,446\,744\,073\,709\,551\,616$. The quotient is printed without padding and the remainder is
padded to 19 digits: \code{"1" + "8446744073709551616"} $= 18\,446\,744\,073\,709\,551\,616$.
\end{example}

These ``linear'' routines are still quadratic overall, but with a small constant. They handle
numbers up to 2048 digits for parsing (\code{BIGMATH\_PARSE\_DC\_THRESHOLD}) and below 1024
digits for printing (\code{BIGMATH\_TOSTR\_DC\_THRESHOLD}), and the leaves of the
divide-and-conquer recursion below.

\section{Large numbers: divide and conquer}

\subsection{Parsing}

Split the digit string in the middle. Parse each half (recursively), then combine:
\[
\text{value} = \text{high} \times 10^{\,\text{digits in low half}} + \text{low} .
\]
Each level costs one big multiplication of about half the total size, and the recursion has
$\log_2 L$ levels: $T(L) = 2\,T(L/2) + \M(L)$, about $\M(L)\log L$ in total.

\begin{figure}[ht]
\centering
\begin{tikzpicture}[level distance=12mm,
  level 1/.style={sibling distance=62mm},
  level 2/.style={sibling distance=31mm},
  every node/.style={box,font=\footnotesize\ttfamily}]
  \node[fill=sBlue!15] {12345678}
    child {node[fill=sBlue!10] {1234}
      child {node[fill=paper] {12}}
      child {node[fill=paper] {34}}
      edge from parent node[draw=none,fill=none,left,font=\footnotesize\rmfamily] {high}}
    child {node[fill=sBlue!10] {5678}
      child {node[fill=paper] {56}}
      child {node[fill=paper] {78}}
      edge from parent node[draw=none,fill=none,right,font=\footnotesize\rmfamily] {low}};
\end{tikzpicture}

\smallskip
{\small Parse upward: $12 \cdot 10^2 + 34 = 1234$, \ $56\cdot 10^2 + 78 = 5678$, \
$1234 \cdot 10^4 + 5678 = 12\,345\,678$.}
\caption{Divide-and-conquer parsing. Leaves use the 18-digit method; each join is one
multiplication by a power of ten.}
\end{figure}

The powers of ten come from \code{Pow10}, which caches them per thread and builds even powers
by \emph{squaring} ($10^{2k} = (10^k)^2$, using the faster square path). After the first
conversion of a given size, the powers are free.

\subsection{Printing}

Printing is the same tree read downward: divide by $10^{k}$, where $k$ is half the digit count.
The quotient gives the top half of the digits, the remainder the bottom half. Recurse on both.

\begin{figure}[ht]
\centering
\begin{tikzpicture}[level distance=12mm,
  level 1/.style={sibling distance=62mm},
  level 2/.style={sibling distance=31mm},
  every node/.style={box,font=\footnotesize\ttfamily}]
  \node[fill=sOrange!15] {10000005}
    child {node[fill=sOrange!10] {1000}
      child {node[fill=paper] {10}}
      child {node[fill=paper] {00}}
      edge from parent node[draw=none,fill=none,left,font=\footnotesize\rmfamily] {$\div 10^4$}}
    child {node[fill=sOrange!10] {0005}
      child {node[fill=paper] {00}}
      child {node[fill=paper] {05}}
      edge from parent node[draw=none,fill=none,right,font=\footnotesize\rmfamily] {$\bmod 10^4$}};
\end{tikzpicture}
\caption{Divide-and-conquer printing of $10\,000\,005$. The remainder $5$ must be written as
\code{0005}: every lower half is padded with zeros to its full width. Only the very top
piece is printed without leading zeros.}
\end{figure}

The divisions are by the same few powers of ten again and again, which is exactly the case
for cached reciprocals (\cref{sec:cached}). For each level of the tree the library stores a
\emph{divider}: the power $10^{k}$ and its Newton reciprocal. The chain of dividers is cached
per thread and reused by later conversions of similar size: the chain's top size is rounded
up to a grid of 16 steps per doubling, so numbers of slightly different lengths share one
chain.

\begin{perfnote}
Before the cached dividers existed, every division in the tree rebuilt its reciprocal, and the
divide-and-conquer printer was \emph{slower} than the quadratic one; it had been tried and
reverted. With cached reciprocals, printing a 100\,000-digit number went from 383\,ms to
45.5\,ms; caching whole divider chains across calls halved warm times again. For small tree
levels (below 13\,000 digits single-threaded, 6\,500 with threads) a cached Newton reciprocal
does not pay, and those levels use the ordinary division dispatcher instead.
\end{perfnote}

\begin{pitfall}
An early version split by $10^{18}, 10^{36}, 10^{72}, \dots$, powers of two times 18 digits,
hoping fewer distinct divisors would help. The splits came out lopsided: a 500\,000-digit
number split into 363\,000 and 137\,000 digits instead of two halves. The recursion lost its
balance and became slower than the quadratic printer (1899\,ms against 1700\,ms). Splits must
be close to the middle; the current rounding keeps them within 56/44.
\end{pitfall}

\subsection{Using several cores}

From 100\,000 digits (\code{BIGMATH\_PARSE\_PARALLEL\_THRESHOLD},
\code{BIGMATH\_TOSTR\_PARALLEL\_THRESHOLD}), both directions descend three levels of the tree on
the calling thread, then hand the $2^3 = 8$ independent subtrees to the thread pool in one
batch. Each worker converts its subtree serially. The calling thread then does the remaining
joins (parsing) or assembles the string (printing). Because every lower half has a known,
fixed width, each subtree can write straight into its slot of the output string.

\section{Cost in practice}

\begin{center}
\small
\begin{tabular}{lrrr}
\toprule
operation & BigMath, 1 thread & GMP, 1 thread & default build $\div$ GMP \\
\midrule
parse 100\,000 digits & 1.644\,ms & 1.057\,ms & 0.52$\times$ \\
parse 20\,000\,000 digits & 910.3\,ms & 809.8\,ms & 0.21$\times$ \\
print 100\,000 digits & 3.186\,ms & 2.382\,ms & 0.52$\times$ \\
print 20\,000\,000 digits & 2\,546\,ms & 2\,118\,ms & 0.24$\times$ \\
\bottomrule
\end{tabular}

\smallskip
{\footnotesize Median warm times from three fresh processes on the M1 Max, revision \code{7b3cc04}. The last column is
time of the default build (up to eight threads) divided by single-thread GMP time.}
\end{center}

\begin{practice}
\begin{itemize}
  \item The first conversion at a new size builds powers of ten and dividers; later ones reuse
        them. Benchmark the second call, and expect the first to be slower.
  \item Workers share immutable CRT transform tables. Each thread reuses writable buffers
        across linear and cyclic products. For full fused matrix products, decimal conversion
        processes the three primes in sequence from length $2^{20}$, using two scratch
        buffers. Standalone multiplication uses this schedule from $2^{24}$.
  \item If no human needs to read the number (storage, hashing, sending it to another
        program), use its limbs directly (\code{x.Limbs()}): that needs no conversion at all.
\end{itemize}
\end{practice}

\newpage
\section{Converting digit vectors between radices}
\label{sec:convertbase}

\code{ConvertBase} takes a little-endian vector of digits or an inclusive
index range within that vector. An equal-radix conversion copies the selected
digits as stored. When the two radices are distinct powers of two, the digits
describe the same bit sequence in different group widths. The implementation
treats \code{Base2\_64}, the radix sentinel with value zero, as a 64-bit width.

The repacker appends each source digit to a 128-bit buffer, starting at the low
end. It emits a target digit whenever enough bits are available, taking the
low bits and shifting them out of the buffer. At most 63 bits are left before
the next source digit arrives. The repacker emits any final partial digit and
trims high zero digits. An empty or all-zero input repacks to \code{\{0\}}.
The work is linear in the combined input and output digit counts.

\begin{example}
In base $2^{31}$, the little-endian pair \code{[0x12345678, 3]} represents
$3\cdot 2^{31} + \texttt{0x12345678}$. In base $2^{16}$ the digits are
\code{[0x5678, 0x9234, 1]}. The high bit of the middle target digit
comes from the low bit of the second source digit.
\end{example}

The remaining radix changes use Horner's rule. Read source digits
from high to low, multiplying the accumulated result by the source radix and
adding the next digit. For a $2^{64}$ source, one step shifts each target digit
left by 64 in a 128-bit word, adds the carry (initially the next source limb),
and divides by the target radix. The remainder replaces the target digit and
the quotient becomes the next carry. See \file{src/algorithms/BaseConversion.cpp}.
Decimal strings use the parsing and formatting routines described above.

\section{Fixed-point decimal arithmetic}
\label{sec:bigdecimal}

A \code{BigDecimal} represents $u\cdot 10^{-s}$, where $u$ is a signed
\code{BigInteger} and $s$ is an integer scale. For example,
\code{12.340} has $u=12340$ and $s=3$. A negative scale represents
trailing integer zeros. Addition and subtraction align scales, and
multiplication multiplies the unscaled integers and adds their scales.
These operations are exact.

\code{a.Divide(b, scale, mode)} returns a value at the requested scale.
If the operands have scales $s_a,s_b$ and the requested scale is $t$,
scale the integer numerator or denominator so the unscaled quotient
represents $(u_a/u_b)\,10^{t+s_b-s_a}$. The integer division produces a
quotient magnitude $q$ and a remainder $r$. For rounding to nearest,
compare $2r$ with the positive divisor. On a tie, \code{HALF\_UP} increases
the magnitude, \code{HALF\_DOWN} leaves it, and \code{HALF\_EVEN} selects
the even result. \code{UP} moves away from zero, \code{DOWN} truncates
toward zero, and \code{CEILING} and \code{FLOOR} round toward the
respective infinities. \code{UNNECESSARY} requires a zero remainder.

\begin{example}
At scale 5, dividing the scaled numerator by 3 gives $100000=3\cdot33333+1$.
\code{BigDecimal("1").Divide(BigDecimal("3"), 5, RoundingMode::HALF\_EVEN)}
returns \code{0.33333}. With \code{UP} it returns \code{0.33334}.
With \code{UNNECESSARY} it throws. A repeating decimal such as $1/3$
is rounded to a finite value at every requested scale.
\end{example}

\code{SetScale} uses the same rounding rules when reducing the scale.
\code{CompareTo} compares numeric values. \code{Equals} also requires
the scales to match, so \code{1.20} and \code{1.2} compare as the same
value while preserving their distinct scales.
